\section*{第 \ref{ch:b3:structural-biology} 章——蛋白质的结构生物学}

\begin{solution}{exo:b3:structural-biology:1}
一级结构：$141$ 个残基的 $\alpha$ 链序列。
\omterm{def:b3:structural-biology:levels}{二级结构}：其八个 $\alpha$ 螺旋。\omterm{def:b3:structural-biology:levels}{三级结构}：
一条链的珠蛋白折叠，包绕其血红素。四级结构：$\alpha_{2}\beta_{2}$
四聚体，及其二聚体彼此相对旋转的方式。
\end{solution}

\begin{solution}{exo:b3:structural-biology:2}
$36\times \qty{0.15}{nm} = \qty{5.4}{nm}$；$36/3.6 = 10$ 圈；
$36 - 4 = 32$ 个氢键。
\end{solution}

\begin{solution}{exo:b3:structural-biology:3}
这说明天然结构（包括四个二硫键在
$105$ 种可能中的正确配对）可自发地从
未折叠的肽链中恢复：序列包含了所需的全部信息，
而天然态是热力学最低点。在仅含纯蛋白的试管中进行，
排除了任何模板、酶或细胞机器
作为信息来源的可能。
\end{solution}

\begin{solution}{exo:b3:structural-biology:4}
NMR：小蛋白，约 \qty{40}{kDa} 以下，在溶液中。
晶体学：任意大小、能结晶的分子，从肽
到核糖体。\omterm{def:b3:structural-biology:methods}{冷冻电子显微镜}：大型复合体，约
\qty{50}{kDa} 以上，无需晶体。NMR 直接反映运动（并且
\omterm{def:b3:structural-biology:methods}{冷冻电子显微镜}还能把分子分入构象类别）；晶体
结构是一种静态平均。
\end{solution}

\begin{solution}{exo:b3:structural-biology:5}
$2^{150} \approx 1.4\times 10^{45}$ 种构象，$1.4\times 10^{33}$
秒 $\approx 5\times 10^{25}$ 年。$3^{20} \approx 3.5\times 10^{9}$，
$3.5$ 毫秒：肽可在一秒内穷举搜索；
而蛋白质在宇宙年龄之内则不能。
\end{solution}

\begin{solution}{exo:b3:structural-biology:6}
$Y = [\alpha(1+\alpha) + Lc\alpha(1+c\alpha)]/[(1+\alpha)^{2} +
L(1+c\alpha)^{2}]$。$\alpha = 1$：$3.01/106 = 0.028$（单一位点
$0.50$）。$\alpha = 10$：$121/242 = 0.50$（$0.91$）。$\alpha = 100$：
$\num{10300}/\num{10601} = 0.97$（$0.99$）。MWC 曲线在低配体浓度下远远
落后于单一位点，随后陡然上升——从 $0.03$
升至 $0.97$，历经两个数量级，而单一位点从 $0.5$ 升至
$0.99$：这种陡峭性即是\omterm{def:b3:structural-biology:allostery}{协同性}。
\end{solution}

\begin{solution}{exo:b3:structural-biology:7}
二磷酸甘油酸仅结合 T 态（位于中央腔中），因此
它稳定 T 态：在模型中它提高 $L$，即未结合配体蛋白的 $T/R$ 比值，
因而使曲线右移——
亲和力降低，组织压强下释放更多氧气。储存的红细胞
失去该化合物，$L$ 下降，曲线左移，输入的
血液保留氧气，而不在组织中释放，直至
细胞在一天内重新合成它。
\end{solution}

\begin{solution}{exo:b3:structural-biology:8}
$d = \lambda/(2\sin\theta) = 0.1/(2\times 0.259) = \qty{0.19}{nm}$：
主链与每个侧链均已定位，单个原子得以分辨。
对于 \qty{0.12}{nm}：$\sin\theta = 0.1/0.24 = 0.417$，$\theta =
\qty{25}{\degree}$。
\end{solution}

\begin{solution}{exo:b3:structural-biology:9}
一个埋于非极性侧链之间的带电赖氨酸，代价是数十
千焦每摩尔（其电荷未被溶剂化，亮氨酸的疏水埋藏
丧失），超过折叠的整个 $\Delta G$：蛋白质
展开，或核心重排。表面上的赖氨酸
与在未折叠态中一样被水合，没有任何损失。
$\qty{-30}{kJ/mol} \approx -12RT$，即 $K = e^{12} \approx 1.6\times
10^{5}$：任一时刻，每 \num{160000} 个分子中有一个处于展开态，而
整个余量相当于数百个氢键中的两三个
——约等于一次突变的效应。
\end{solution}

\begin{solution}{exo:b3:structural-biology:10}
$R$ 态的总量为 $[R_{0}](1+\alpha)^{n}$，$T$ 态为
$[R_{0}]L(1+c\alpha)^{n}$，因此 $\bar R = (1+\alpha)^{n}/[(1+\alpha)^{n}
+ L(1+c\alpha)^{n}]$。当 $c \to 0$ 时 $T$ 项为 $L$，且 $\bar R =
1/2$，当 $(1+\alpha)^{n} = L$ 时，即 $\alpha = L^{1/n} - 1$。对于
血红蛋白 $(3\times 10^{5})^{1/4} = 23.4$，故 $\alpha \approx 22$：
\qty{22}{mmHg}，接近 \qty{26}{mmHg} 的 $P_{50}$——
构象开关与半饱和几乎重合。
\end{solution}

\begin{solution}{exo:b3:structural-biology:11}
此说曾遭抵制，因为所有已知的感染因子都通过核酸进行复制，
而蛋白质无法复制其序列。若
发现必需某种核酸，或纯化蛋白被证明不能
传播，则将推翻它。支持证据是：感染性在核酸酶、
紫外线及足以破坏任何基因组剂量的电离辐射下仍然保留，
但会被蛋白变性剂和蛋白酶破坏；重组
PrP 在体外折叠成纤维后可向动物传播疾病，
且传代后保持相同的株特性。需要种子，是因为
正常单体的转化单独进行时极为缓慢——
交叉 $\beta$ 结构的核必须先形成，这是
限速事件——而已形成的纤维末端
可迅速转化单体：形状通过模板化生长而被复制。
\end{solution}

\begin{solution}{exo:b3:structural-biology:12}
许多序列折叠为同一结构：折叠所需的是一种
疏水与极性位置的排布模式，少数甘氨酸和脯氨酸
位于转角处，以及执行功能的残基；表面上的其余部分
可无代价地改变，因此突变在那里积累，而
选择则保留折叠与功能。保守残基因而是
核心（与其说是具体残基，不如说是一种模式）、
催化与结合残基，以及结构性的甘氨酸、脯氨酸和
半胱氨酸。谱方法正是对这些列加权，因此
能在 \qty{15}{\%} 的序列一致性下识别家族成员；蛋白质设计则
逆向运用这一逻辑，为目标折叠选择疏水模式，
而表面则可任意设置。
\end{solution}

\begin{solution}{pb:b3:structural-biology:1}
\textbf{1.} $0.75\times 150 \approx 112$ 个螺旋残基；$112\times
\qty{0.15}{nm} = \qty{17}{nm}$；$112/3.6 \approx 31$ 圈。
\textbf{2.} $112 - 8\times 4 = 80$ 个氢键。
\textbf{3.} 质量 $150\times 110 = \qty{16500}{Da} = 2.7\times 10^{-20}$
克；体积 $2.7\times 10^{-20}\times 0.73 = 2.0\times 10^{-20}$ cm$^{3}$
$= \qty{20}{nm^3}$；半径 $(3V/4\pi)^{1/3} = \qty{1.7}{nm}$。
\textbf{4.} 伸展长度 $150\times 0.36 = \qty{54}{nm}$，与
直径 \qty{3.4}{nm} 相比，相差 $16$ 倍。
\textbf{5.} 螺旋残基相当于伸展长度的 \qty{75}{\%}，
即 $112\times 0.36 = \qty{40}{nm}$；螺旋将其压缩至
\qty{17}{nm}，压缩倍数为 $0.36/0.15 = 2.4$。
\textbf{6.} 获得 $40\times 4 = \qty{160}{kJ/mol}$，失去 $150\times 0.9 =
\qty{135}{kJ/mol}$；净 $\Delta G \approx \qty{-25}{kJ/mol}$，
处于临界范围。
\textbf{7.} $3^{150} \approx 4\times 10^{71}$；$4\times 10^{59}$ 秒
$\approx 10^{52}$ 年。
\textbf{8.} $10^{-2}\times 10^{12} = 10^{10}$ 种构象，
占总数的 $3\times 10^{-62}$。
\textbf{9.} 各螺旋段在 \qty{1}{\micro s} 内完成（并行）；$8! = \num{40320}$ 种
堆积方式以每秒 $10^{6}$ 种的速度检验，需 \qty{40}{ms}：总计约
\qty{40}{ms}，与观测到的数量级一致。是的：先独立求解各部分
再求整体，这正是一个\omterm{thm:b3:structural-biology:levinthal}{折叠漏斗}——搜索空间
被分解，而非被穷尽。
\textbf{10.} 期望循环数 $1/0.3 = 3.3$，约 \qty{33}{s}；
六次循环后，$0.7^{6} = 0.12$ 仍未折叠。
\textbf{11.} $\Delta G = -RT\ln K$，其中 $K = 10^{4} \to 10^{2}$：由
$-23.7$ 变为 \qty{-11.9}{kJ/mol}，稳定性降低 \qty{12}{kJ/mol}
（$RT = \qty{2.58}{kJ/mol}$，$\ln 100 = 4.6$）。
\textbf{12.} 聚集始于一个核的形成，其形成速率
随易于聚集的未折叠物种浓度的高次幂而上升；该物种浓度
增加一百倍，可使种子在寿命期内形成的概率
提高数个数量级，一旦
出现种子，生长迅速，并能自我模板化。
\textbf{13.} 当 $\alpha = 100$ 时：$Y = 1.054\times 10^{8}/1.089\times 10^{8} =
0.97$。当 $\alpha = 20$ 时：分子 $185\,220 + 103\,680 = 288\,900$，
分母 $194\,481 + 622\,080 = 816\,561$：$Y = 0.35$。
\textbf{14.} 在肺中 $\bar R = 1.041\times 10^{8}/1.089\times 10^{8} = 0.96$；
在肌肉中为 $194\,481/816\,561 = 0.24$。
\textbf{15.} 卸载的容量为 $0.97 - 0.35 = 0.62$（约占
负荷的三分之二）。单一位点：$100/126 - 20/46 = 0.79 - 0.43 = 0.36$。
\textbf{16.} 容量 $150\times 1.34 = \qty{200}{mL/L}$；输往肌肉的
量为 $0.62\times 200 \approx \qty{125}{mL/L}$。静息时 $(0.97 - 0.78)\times
200 = \qty{38}{mL/L}$；$250/38 \approx \qty{6.5}{L/min}$ 的
血流量，与静息心输出量的数量级相当。
\textbf{17.} 当 $L = 3\times 10^{6}$，$\alpha = 40$ 时：分子 $2.76\times
10^{6} + 3.29\times 10^{6} = 6.05\times 10^{6}$，分母 $2.83\times
10^{6} + 11.5\times 10^{6} = 14.4\times 10^{6}$：$Y = 0.42$，而非
$0.78$——静息时卸载的氧显著增加，体现了血液对高海拔的适应
与肌肉对工作负荷的适应。
\textbf{18.} 胎儿血：$L = 3\times 10^{4}$，$\alpha = 30$，分子
$893\,730 + 19\,770 = 913\,500$，分母 $923\,520 + 85\,680 =
1\,009\,200$：$Y = 0.91$。母体在相同压强下，分子
$893\,730 + 197\,730 = 1\,091\,460$，分母 $923\,520 + 856\,830
= 1\,780\,350$：$Y = 0.61$。在相同压强下，胎儿血液结合的
氧更多，因此氧经胎盘由母体流向胎儿。
\textbf{19.} $(10^{5}\,\text{nm}/6\,\text{nm})^{3} = 4.6\times 10^{12}$
个细胞；$1.9\times 10^{13}$ 个分子。
\textbf{20.} $d = 0.1/(2\sin 14.5^{\circ}) = \qty{0.20}{nm}$：是的，
所有侧链与原子均已定位。
\textbf{21.} $\sin\theta = 0.1/0.3$，$\theta = \qty{19.5}{\degree}$。在
无序晶体中，分子排列并非严格对准，
在大角度处散射的波（携带精细细节信息）
不再同相叠加，大角度斑点消退，只有低\omterm{def:b3:structural-biology:methods}{分辨率}
的斑点保留下来。
\textbf{22.} 将 $d$ 减半需要两倍的信号，因此需要四倍的
颗粒数：约 \num{40000}（实际需要更多，因其他误差会引入）。
\textbf{23.} 膜蛋白，它们很少能从去污剂
溶液中结晶，但可在脂质纳米盘中成像；大型柔性
复合体——如核糖体、剪接体，以及带有调节因子的离子通道，
其异质性使其难以结晶，但其颗粒
可按构象分入不同类别。
\textbf{24.} 结合的血红素与氧及其精确几何构型、
有序排列的水与离子、第二种构象（T 态对 R 态）以及
折叠正确的证明——预测仅是一个未结合配体的
模型，尚无实验检验。
\textbf{25.} 卸载容量的 $0.62$（单一位点为 $0.36$）；
莱文塔尔时间约 $10^{52}$ 年；\omterm{def:b3:structural-biology:methods}{分辨率}达 \qty{0.20}{nm} 的图。
\end{solution}
